Let \(E\) erase only the ordered provenance list \(\Lambda\) from a trace-aware state.  Write the initial state as
\[
\kappa=(\mathcal H,F,c,q,\pi,\Lambda).
\tag{1}
\]
The premise that both colored composites are legal supplies fixability at each intermediate CADMG, validity of both licenses, and support of both denominator traces.  In addition, \(\text{ass:strict-positivity}\) makes every conditional kernel appearing in a legal fixing strictly positive on its stated support.  Thus every division below is defined.

Apply \(\text{lem:ordinary-adjacent-fixing-commutation}\) to the CADMG \(\mathcal H\) and normalized kernel \(q\).  Its hypotheses are the two fixability clauses in \(\operatorname{Swap}_{\kappa}(u,v)\), together with the assumed legality of both composites.  It gives
\[
\left(E\Phi_v^{\lambda_v}\Phi_u^{\lambda_u}(\kappa)\right)_{\mathcal H,q}
=
\left(E\Phi_u^{\lambda_u}\Phi_v^{\lambda_v}(\kappa)\right)_{\mathcal H,q}.
\tag{2}
\]
In either order the fixed set is \(F\cup\{u,v\}\), because each fixing adds its vertex and never removes a fixed vertex.  Deleting arrowheads into \(u\) and into \(v\) commutes, and fixing does not alter the context assignment \(c\).  Hence the CADMG, fixed-set, and context components in the two erased states agree.

By \(\text{def:colored-fixing}\), the provenance expression \(\pi\) is transformed by exactly the same two supported conditional-kernel divisions as \(q\).  Equality (2), viewed as the equality obtained after those two divisions, therefore gives equality of the resulting rational observed-law functions after cancellation of their nonzero denominators.  Thus
\[
E\Phi_v^{\lambda_v}\Phi_u^{\lambda_u}(\kappa)
=
E\Phi_u^{\lambda_u}\Phi_v^{\lambda_v}(\kappa)
\tag{3}
\]
in every component retained by \(E\).

It remains to calculate the component erased by \(E\).  Starting from the provenance \(\pi\), the order \(u,v\) appends
\[
\bigl((\lambda_u,\pi),(\lambda_v,\Phi_u^{\lambda_u}\pi)\bigr),
\tag{4}
\]
whereas the order \(v,u\) appends
\[
\bigl((\lambda_v,\pi),(\lambda_u,\Phi_v^{\lambda_v}\pi)\bigr).
\tag{5}
\]
Equations (4)--(5) follow directly from the append rule in \(\text{def:colored-fixing}\).  Merely permuting the two list positions does not replace \(\Phi_u^{\lambda_u}\pi\) by \(\pi\), or \(\Phi_v^{\lambda_v}\pi\) by \(\pi\); hence the hypotheses do not imply equality of the full ordered trace states up to permutation.  Finally, both branches have the same pre-swap state (1) and the same two licenses.  Replaying those licenses from (1) in a fixed canonical order deterministically recomputes one of (4)--(5), while (3) proves equality of all non-list components.  Therefore a branch merge is sound only after that canonical replay, as claimed.